\chapter{Background and Conceptual Basis}\label{ch:literature}

\section{Introduction}\label{sec:ch2_introduction}

This chapter reviews how transmission, viral evolution, and surveillance practices shape genomic evidence of recent linkage and possible superspreading.

We begin with incomplete sampling and the temporal and lineage contexts for comparing genomes, before reviewing methods for pairwise inference and clustering. We then examine the basis for screening cluster patterns through time and characterising the populations and surveillance contexts they represent.

\section{Genomic Surveillance and Incomplete Sampling}\label{sec:ch2_partial_observation}

Within public health surveillance, genomic data provide additional evidence about pathogen circulation and transmission. This section examines how surveillance objectives and sampling practices shape that evidence, using the UK COVID-19 response to illustrate the implications of incomplete observation.

\subsection{Surveillance as partial observation}\label{sec:ch2_partial_obs_concept}

Public health surveillance encompasses the continual collection, analysis, interpretation, and use of health-related data to guide public health responses. For infectious diseases, these activities help determine where transmission occurs and how it is sustained \citep{Murray2017,Groseclose2017}. Genomic surveillance extends this framework by integrating pathogen sequence data with diagnostic, epidemiological, clinical, and contextual information \citep{WHO2023,Tiwari2025,ECDC2014,Black2020}. Genomes can help identify possible links between cases with no previously recognised epidemiological connection. Associated metadata provide the information about people, places, and dates needed to interpret these genetic relationships \citep{Gardy2017,Armstrong2019,Carter2022}.

For an infection to appear in a genomic surveillance dataset, a sample must be collected and tested, its viral genome successfully sequenced and pass quality control, and the result reported and linked to the required metadata. Inclusion therefore depends on testing availability and eligibility, test-seeking behaviour, laboratory capacity, sequencing throughput, and quality-control choices, while turnaround time, reporting delays, and metadata completeness affect when and how the resulting records can be used. These conditions vary across time, geography, clinical settings, laboratories, and social groups, producing a biased and often delayed sample of the epidemic \citep{Chen2022,Han2023,Brito2022}. Genomic surveillance is therefore a \emph{partial observation system}: it captures only those infections that pass through these steps.

\subsection{Surveillance priorities shape sampling}\label{sec:ch2_surveillance_priorities}

The purpose of surveillance helps determine which infections are selected for sequencing.\ \emph{Variant surveillance} aims to detect and track lineages exhibiting changes in transmissibility, severity, immune escape, or diagnostic and therapeutic performance \citep{WHO2021,ECDC2021,OudeMunnink2021}.\ \emph{Localised outbreak investigation} uses genomic and epidemiological data to assess transmission in predefined settings such as hospital wards, care homes, schools, workplaces, prisons, or other institutions \citep{Meredith2020,Hamilton2021,Aggarwal2022,Sekizuka2020,Reichert2022,Adamson2022}.\ \emph{Population surveillance} describes viral introductions, changes in circulating lineages by region, and the changing composition of the epidemic over time \citep{daSilvaFilipe2020,duPlessis2021,Lycett2021,Vohringer2021}. The same genomic record can contribute to all three aims, but each places different demands on sampling.

Broad, representative sampling supports population surveillance, while variant surveillance may use sparse but broad coverage alongside targeted sequencing of travel-associated or suspected variant cases. Outbreak investigation may require dense sampling within a setting, including sequencing for hospital or care-home infection control \citep{WHO2021,ECDC2021,OudeMunnink2021,Inward2022}. Who is observed therefore reflects infection risk as well as the sampling strategy used. These strategies help answer public health questions with limited resources, but datasets assembled for different purposes are not automatically comparable. Low testing rates can reduce surveillance sensitivity, while uneven sampling across place, time, and investigation type can bias variant detection, lineage prevalence estimates, and inferred cluster size or composition \citep{Suster2022,Han2023}. Recent-linkage inference in this thesis must consequently consider both how infections were sampled and whether the comparison genomes come from a relevant period and evolutionary background.

\subsection{Scale and completeness in UK genomic surveillance}\label{sec:ch2_uk_surveillance_scale}

The COVID-19 response drew on established sequencing platforms and data infrastructure. These included open data-sharing systems such as the Global Initiative on Sharing All Influenza Data (GISAID)\nomenclature{GISAID}{Global Initiative on Sharing All Influenza Data} \citep{Khare2021} and phylogenetic visualisation systems like Nextstrain \citep{Hadfield2018}. In the United Kingdom (UK)\nomenclature{UK}{United Kingdom}, the COVID-19 Genomics UK Consortium (COG-UK)\nomenclature{COG-UK}{COVID-19 Genomics UK Consortium} developed a national sequencing network linking public health agencies, universities, sequencing laboratories, and data infrastructure \citep{COGUK2020,Nicholls2021}. This infrastructure enabled large-scale analyses of viral introductions, onward transmission, lineage establishment, variant replacement, and spatial spread, including studies focused on Scotland and the wider United Kingdom \citep{daSilvaFilipe2020,duPlessis2021,Lycett2021,Vohringer2021}.

Scale alone, however, does not ensure complete observation. During the COVID-19 pandemic, sequencing in Scotland and the wider United Kingdom depended on changing testing pathways, variant-monitoring priorities, outbreak investigations, laboratory capacity, policy, and epidemic waves. These surveillance conditions limit what the sampled data can establish about recent transmission and possible superspreading.

As outlined in Chapter~\ref{ch:introduction} (Section~\ref{sec:ch1_genomic_surveillance}), pathogen genomes do not provide a direct record of transmission pathways. Dense sampling within an area can make compatibility based on genetic differences and sampling dates informative about recent linkage. With sparse, uneven, or targeted sampling, apparent clusters or gaps can also reflect which infections were detected and sequenced. Where a lineage is highly prevalent, genetic similarity may reflect background circulation rather than recent shared exposure. Genomic clusters should therefore be interpreted as the joint product of pathogen evolution, transmission, and surveillance. Chapter~\ref{ch:genomic_networks} returns to this point empirically: it describes how the numbers of genomes sequenced, sequencing coverage, policy context, and circulating clades changed over the period analysed in this thesis.

\section{Temporal and Lineage Comparison Contexts}\label{sec:ch2_windows_lineages}

Comparing every pair of genomes in the national dataset includes many pairs too distant in time or genetics to be informative about recent transmission. This section explains why comparisons are restricted by sampling period and viral lineage.

\subsection{Bounding comparisons in time}\label{sec:ch2_temporal_scales}

SARS-CoV-2 causes an acute respiratory infection in which onward transmission is concentrated over days to weeks, whereas the national genomic record spans years. The relevant timescale depends on the incubation period, serial interval, generation time, infectiousness profile, and how viral shedding changes over the course of infection (Figure~\ref{fig:disease_timeline}). Early studies estimated incubation and serial intervals at approximately five to six days, with generation times typically ranging from four to six days \citep{McAloon2020,Alene2021,Nishiura2020,Du2020,Ganyani2020,Ferretti2020,Hart2022,Kim2023}. These parameters vary by variant, with shorter generation times generally reported for Omicron than for earlier lineages during the COVID-19 pandemic \citep{Xu2023,Madewell2023}. The critical implication is that these estimates support assessing recent linkage over days to weeks, but do not establish a single temporal cutoff. The timing assumptions used for inference must therefore accommodate uncertainty in generation-time estimates, particularly when a new variant emerges and data are limited.

\begin{figure}[htbp]
    \centering \includegraphics[width=\linewidth]{infectious_disease_timeline} 
    \caption[Time periods and intervals in infectious-disease transmission]{\textbf{Time periods and intervals in infectious-disease transmission.} Schematic timeline for a single infector-infectee pair; ordering and relative durations are illustrative and vary by pathogen, host, and context.\ \emph{Within-host} periods are defined relative to an individual's infection: the \emph{latent period} (infection to onset of infectiousness), the \emph{infectious period} (window permitting onward transmission), and the \emph{incubation period} (infection to symptom onset).\ \emph{Between-host} intervals connect the pair: the \emph{generation time} separates infection times, whereas the \emph{serial interval} separates symptom-onset times. The generation time is strictly positive. The serial interval is negative when the infectee develops symptoms before the infector.}\label{fig:disease_timeline}
\end{figure}

This short timescale for onward transmission contrasts with the slower pace of changes visible in SARS-CoV-2 consensus genomes. The virus evolves at approximately two substitutions per genome per month, meaning that many transmission events can occur before a new difference appears in the sampled consensus genomes \citep{Duchene2020,Pybus2013}. Detectable genetic change therefore accumulates more slowly than transmission occurs. Genetic identity between genomes sampled far apart in time does not by itself establish recent transmission between hosts. Two identical sequences collected months apart may be connected through a long chain of unsampled intermediates, represent independent descendants of a widely circulating viral genetic sequence, or result from convergent evolution \citep{Martin2021}.

Similar viral genomes can continue to be sampled over time, providing another reason to restrict comparisons by collection date. In a growing epidemic, sampled lineages tend to be concentrated around the period of expansion, but descendants of the same ancestral viral sequence can be sampled later or in different geographic and epidemiological settings \citep{Volz2009,Volz2013}. Earlier samples can therefore remain genetically similar to later samples long after the main epidemic wave, especially when viral evolution is slow or sampling is sparse. Such persistence may arise from low-level circulation, repeated reintroduction, undetected transmission, or prolonged infections. Prolonged infections can also generate unusual mutation patterns as the virus continues to evolve within one person \citep{Choi2020,Kemp2021}. When samples from different periods are pooled, these shared genetic patterns can link clusters even when the data do not support recent transmission between them.

Restricting comparisons to rolling time windows keeps each analysis focused on samples collected within a recent period. The analysis window should span relevant infection-to-sampling intervals and several transmission generations, while remaining short enough to limit comparisons with older samples whose genetic similarity may obscure recent patterns. Reporting delays separately affect which sequences are available when the analysis is run.

\subsection{Bounding comparisons by lineage}\label{sec:ch2_co_circulating_subepidemics}

A national SARS-CoV-2 epidemic can be viewed as a set of partly overlapping sub-epidemics with different transmission histories and viral ancestries \citep{Chowell2019,Vohringer2021}. Phylogenetic studies support this view, showing that Scotland's epidemic involved multiple introductions followed by local expansion of specific viral lineages \citep{daSilvaFilipe2020,duPlessis2021,Lycett2021}.

Lineage assignment systems such as Pango formalise this structure by assigning dynamically updated labels to genetically coherent and epidemiologically meaningful groups of SARS-CoV-2 viruses \citep{Rambaut2020,OToole2021}. These labels do not themselves define transmission clusters, but provide evolutionary context for recent-linkage inference. Viruses assigned to distinct major lineages typically differ by substitutions accumulated over many transmission generations and have different transmission histories. Large genetic distances therefore provide little support for recent linkage under the assumed evolutionary model. Conversely, small distances can be misleading for viruses whose ancestry is uncertain or affected by recombination.

Grouping comparisons by lineage also accounts for differences in transmissibility, interactions with population immunity, and epidemic growth. Alpha, Delta, and Omicron lineages showed different growth advantages or abilities to evade immunity \citep{Volz2021,Davies2021,Mlcochova2021,Viana2022,Cao2021}. Viral genetic variation is shaped by mutation rates, selection pressures, founder effects, and transmission bottlenecks. The differences observed between samples also depend on which infections were sequenced, the number of unsampled intermediates, and sequencing error \citep{Duchene2020,Ghafari2022,Pybus2013}. Combining epidemic periods and lineages without accounting for their differences can make equal-sized clusters difficult to compare epidemiologically.

\subsection{A scalable comparison design}\label{sec:ch2_scalable_design}

The preceding arguments support using both collection dates and lineage assignments when assessing recent linkage. Collection dates help determine whether genetic similarity is informative about recent transmission. A shared lineage label provides evolutionary context, but neither close collection dates nor a shared lineage label alone establishes recent linkage.

Dividing the national surveillance record by sampling period and lineage produces smaller comparison groups. Within each group, genetic and collection-date differences, and the clusters inferred from them, can be assessed within a shared lineage background and a recent sampling interval. These groups make the comparisons easier to interpret and reduce the number of sequence pairs that must be analysed. Comparing every pair among hundreds of thousands of genomes is computationally expensive and includes many pairs that provide little information about recent transmission.

Existing SARS-CoV-2 genomic systems make large datasets manageable in different ways: subsampling and continuous updating in Nextstrain \citep{Hadfield2018}, maximum-parsimony placement onto mutation-annotated trees in Ultrafast Sample placement on Existing tRees (UShER)\nomenclature{UShER}{Ultrafast Sample placement on Existing tRees} \citep{Turakhia2021}, and hierarchical lineage summarisation in Pango \citep{Rambaut2020,OToole2021}. Similarly, molecular cluster-detection approaches use thresholds, partitions, or precomputed trees to identify groups of sequences for investigation without evaluating every possible transmission pair \citep{Prosperi2011,Balaban2019,Gamza2024}.

Grouping sequences by time and lineage defines comparison sets for assessing recent-transmission compatibility. It makes those comparisons more relevant to the chosen target, although sampling differences and variation within each group still affect interpretation.

\section{Phylodynamics and Pairwise Inference}\label{sec:ch2_common_approaches}

Phylodynamics uses pathogen genealogies to study epidemic processes, including growth, introductions, migration, and variant replacement \citep{Grenfell2004,Pybus2013,Volz2013,Lemey2009,Faria2011,Suchard2018}. These analyses describe ancestry and population patterns, but do not automatically identify recent transmission between individual infections. A phylogenetic tree describes the ancestry of sampled pathogen sequences; a transmission tree describes who infected whom. Within-host diversity, transmission bottlenecks, and incomplete sampling mean that a branching event in one tree need not correspond to a transmission event in the other \citep{Ypma2013,Jombart2014,DeMaio2015,Didelot2017}. For SARS-CoV-2, consensus genomes can also remain identical across epidemiologically distinct infections \citep{Valesano2021,Lythgoe2021,Bendall2022,Farjo2023}. The same small genetic distance may therefore be consistent with direct transmission, an unsampled chain, or infection from a shared source \citep{Louca2021}.

Methods for reconstructing transmission address this distinction. Outbreaker2 combines genetic and temporal data to infer transmission histories using Bayesian methods \citep{Campbell2018b,Jombart2014}. TransPhylo uses a dated phylogeny to infer transmission while allowing for within-host diversity, unsampled cases, and an outbreak that is still ongoing \citep{Didelot2017}. SCOTTI models pathogen evolution within hosts and transmission between them, including the possible involvement of unsampled hosts \citep{DeMaio2016}. These methods quantify uncertainty about transmission histories under their respective models. Their interpretation depends on the available genomic and epidemiological information and on assumptions about transmission, evolution, and sampling.

Detailed reconstruction is valuable when an investigation seeks likely infectors or transmission routes within a bounded outbreaks where the case population, sampling process, and epidemiological context are reasonably well characterised \citep{Meredith2020,Abbas2022,Duault2022,Campbell2018b}. The question addressed by routine surveillance may be narrower: do two sampled infections have genetic and temporal differences consistent with recent linkage? This question does not require a complete transmission tree. It does require a clear definition of the relationship being assessed and an account of the uncertainty that remains.

Tree-based clustering groups sequences using their positions and distances on an inferred phylogeny. Cluster Picker combines genetic-distance and branch-support criteria, while TreeCluster offers several ways to partition a tree using branch-length constraints \citep{RagonnetCronin2013,Balaban2019}. These methods use shared evolutionary structure without reconstructing each transmission event. Their results depend on the inferred tree, the sampled sequences, and the clustering rule. A tree-derived cluster can help identify an outbreak-associated group, but membership alone does not establish a particular transmission relationship \citep{Bousali2021,Lemieux2021,GomezCarballa2021}.

Fixed genetic thresholds offer a simpler pairwise rule: two sequences are linked if their genetic distance is below a chosen value. Such rules are widely used in bacterial and human immunodeficiency virus (HIV)\nomenclature{HIV}{human immunodeficiency virus} surveillance and are straightforward to apply consistently \citep{Maiden1998,KosakovskyPond2018}. Their epidemiological meaning depends on the pathogen and surveillance setting. For a low-diversity pathogen such as SARS-CoV-2, the same distance can occur across different transmission histories, sampling intervals, and lineage backgrounds \citep{Poon2016,Sobkowiak2022}. Different studies have consequently used different thresholds \citep{Hare2023}. A genetic cut-off defines which pairs are retained, but does not by itself explain why those pairs should represent recent linkage.

Pairwise models can incorporate sampling dates alongside genetic distance. For example, cov2clusters estimates the probability of shared cluster membership using a logistic model with patristic distance (the summed branch lengths between two tips) and differences in sampling or symptom-onset dates \citep{Sobkowiak2022}. Its coefficients can be set manually or fitted using known cluster memberships. Fitting requires suitable labels, while manual specification requires justification for the chosen values. Because cluster membership need not imply direct transmission, these probabilities must be interpreted in relation to the clustering target, how the coefficients were obtained, and the population studied \citep{Sobkowiak2022,Sobkowiak2023}.

Pairwise methods also differ in the transmission relationships they assess. A2B-COVID evaluates whether observations from two individuals are consistent with direct transmission from one to the other \citep{Illingworth2022}. This target differs from assessing recent linkage that also includes co-primary infection from a shared source. Transcluster estimates the number of intermediate transmissions separating sampled infections using genetic differences, sampling dates, and assumed evolutionary and transmission rates \citep{Stimson2019}. Although illustrated using tuberculosis data, it has also been applied to SARS-CoV-2 \citep{GallegoGarcia2022}.

Infection dates are usually unknown, and variable delays to testing or sampling affect the observed interval between cases. Mutation accumulation also varies, so the same type of transmission relationship can produce different genetic distances \citep{Jombart2014,Didelot2017,DeMaio2016}. The variability represented by such models should be distinguished from uncertainty about their parameters and assumptions.

The pairwise measure required here must therefore combine genetic differences and sampling dates to assess explicitly defined recent-linkage scenarios without fitting to labelled transmission pairs. Chapter~\ref{ch:epilink} develops and evaluates EpiLink for this purpose, with the primary analyses assessing direct transmission and co-primary infection from a shared source under explicit model assumptions. Its scored scenarios can include additional unsampled intermediates when the investigation targets longer transmission chains.

\section{Pairwise Linkage and Genomic Clusters}\label{sec:ch2_compatibility_graphs}

Once pairwise relationships between sampled infections are estimated, they can be represented as a weighted or unweighted graph \citep{Liu2023,KosakovskyPond2018,Stimson2019,Sobkowiak2022,Illingworth2022}. Each sampled genome corresponds to a node, and each retained pairwise relationship forms an edge. In a weighted graph, the edge weight describes the strength of that relationship (Figure~\ref{fig:pairwise_clustering_schematic}A). Removing edges with very low weights, a step called graph sparsification, can focus the analysis on stronger relationships and reduce the computational cost of analysing graphs with many connections. In this thesis, edge weights describe compatibility with selected recent-linkage scenarios. The resulting weighted compatibility graph provides the basis for the graph and cluster analyses.

The meaning of a cluster depends first on what those edges represent. Genetic similarity, a predicted probability of shared cluster membership, and compatibility with specified transmission scenarios describe different kinds of relationship. Community detection can identify groups within graphs based on any of these measures, but cannot turn an uncertain pairwise relationship into a confirmed transmission link. Interpreting a cluster therefore requires both a clear definition of its edges and an explanation of how those edges are grouped \citep{Liu2023,Poon2016}.

\begin{figure}[htbp]
    \centering
    \includegraphics[width=\textwidth]{pairwise_clustering_schematic}
    \caption[Weighted compatibility graphs and cluster definitions]{\textbf{Weighted compatibility graphs and cluster definitions.} (A) Pairwise compatibility scores define a weighted graph. Thick solid edges indicate stronger compatibility; pale solid edges indicate weaker retained compatibility; dotted edges represent very weak candidate relationships that may be removed during sparsification. (B) Weighted community detection partitions the full graph according to stronger compatibility within groups. (C) After removing very weak edges, the main communities remain visible in this example. However, retained weak bridges join them into the same connected component when clusters are defined as groups joined by paths of retained edges. (D) In this example, weighted community detection retains the main groups after sparsification because weak bridges carry less weight than the many strong connections within each group.
    }\label{fig:pairwise_clustering_schematic}
\end{figure}

Information about cases and samples can be attached to graph nodes and edges. For instance, collection date, geography, age, sex, deprivation index, vaccination status, policy period, testing context, and other surveillance metadata can be attached to nodes. Genetic distance, collection-date difference, compatibility score, and uncertainty summaries can be assigned to edges. Variables used to construct the graph, such as collection dates and genetic distances, have a different role from variables added to describe it, such as demographic attributes. The same variable can serve both roles.

Population composition and the pattern of connections answer different questions. Composition describes which groups contribute sampled genomes. For example, two graphs could contain the same proportions of two age groups, with the same total connection weight attached to each group. In one graph, most weight could connect members of the same age group; in the other, weight could be distributed more evenly between connections within and between age groups. The composition would be identical, but the organisation of compatibility would differ.

Weighted categorical assortativity summarises this organisation relative to a reference based on the observed graph \citep{Newman2002,Newman2003}. Node strength is the sum of the weights attached to a sampled genome. Each category therefore accounts for a share of total node strength. The reference asks how much compatibility weight would connect the same category if edge endpoints were paired at random while preserving these category strength shares. Positive assortativity indicates more within-category weight than this reference, and negative assortativity indicates less. This is a comparison of compatibility weights, not an assumption that people make contacts at random. The reference also differs from a comparison based only on the number of sampled genomes in each category.

Contact research provides an epidemiological reason to examine such patterns. Network models show that the organisation of contacts can affect epidemic dynamics and the performance of contact tracing \citep{Kiss2007}. Empirical contact surveys found substantial age-assortative contact, particularly among schoolchildren and young adults \citep{Mossong2008}. These findings motivate asking whether genomic compatibility also has age structure. They do not supply an expected coefficient for a compatibility graph, whose edges represent genetic and temporal evidence rather than observed contacts. Likewise, genomic studies of introductions and subsequent circulation in Scotland and the wider United Kingdom motivate examining whether compatibility is concentrated within residential regions \citep{daSilvaFilipe2020,duPlessis2021}. Geographic differences in testing and sequencing can contribute to the same pattern.

The attributes describe different aspects of population structure. Health board and local authority identify administrative areas of residence; urban/rural class identifies a type of settlement, which can be shared by people living far apart. The Scottish Index of Multiple Deprivation (SIMD) groups areas by deprivation, so membership of the same quintile also need not imply proximity. Differences in infection risk, severe disease, or ascertainment can change the representation of demographic and deprivation groups without establishing whether compatibility weight preferentially connects members of the same group. There is therefore no general expectation that sex or SIMD assortativity must be positive, or that the ordering of coefficients measures the relative importance of these attributes for transmission.

In the Scottish application, age, sex, deprivation, and residence are used to describe graphs whose pairwise scores are calculated from genetic differences and collection dates. Their associations with compatibility provide additional epidemiological context, conditional on lineage, sampling period, and which infections were observed. Chapter~\ref{ch:genomic_networks} uses these literature-informed expectations to interpret exploratory comparisons across attributes, windows, and lineages. Agreement with an expected population pattern can support its epidemiological plausibility, but cannot independently validate individual transmission links.

Other summaries describe the connections themselves. Degree counts the retained links attached to a genome, while strength sums their weights \citep{Newman2003}. Degree and strength assortativity ask whether links preferentially join genomes with similar numbers of connections or similar total connection weights. These describe the structure of the sampled compatibility graph; their interpretation differs from the numbers of contacts or secondary infections associated with a person. Edge-level summaries can additionally describe compatibility weights, genetic distances, and collection-date differences among retained pairs.

Clusters summarise groups of sampled genomes connected within the weighted compatibility graph. The simplest approach removes edges below a specified threshold and identifies connected components as clusters (Figure~\ref{fig:pairwise_clustering_schematic}C). Connected components are transparent, scalable, reproducible, and easily communicated, explaining their widespread use in pairwise genomic clustering \citep{KosakovskyPond2018,Stimson2019,Sobkowiak2022,Sobkowiak2023}. However, this simplicity has limitations \citep{Liu2023,Poon2016}. A small change in the threshold can retain or remove a bridge joining large groups. The method assigns samples to the same cluster whenever a path of retained edges connects them, regardless of the weights along that path. Edge weights can still be analysed separately, but do not affect component membership once the threshold has been applied.

Community detection offers an alternative approach to summarising weighted compatibility graphs. Rather than assigning all nodes joined by paths to a single group, community-detection algorithms seek partitions in which nodes are more strongly connected within groups than between them \citep{Liu2023,Fortunato2010,Newman2004}. Edge weights distinguish stronger from weaker relationships. A weak bridge need not force two groups into a single cluster when connections within each group are much stronger than the connection between them (Figure~\ref{fig:pairwise_clustering_schematic}B and D).

Several broad families of methods are available for this purpose. Modularity-based methods compare observed within-community connections with those expected under a null graph that accounts for node degrees or strengths \citep{Newman2004,Newman2006}; the Louvain algorithm made this optimisation practical at large scales \citep{Blondel2008}. Flow-based methods, stochastic block models, spectral methods, and hierarchical approaches use different definitions of what makes a group of nodes a community \citep{Rosvall2008,Holland1983,Fortunato2010}. These methods can therefore divide the same graph into different groups.

Leiden is one algorithm for finding communities in weighted graphs. It extends the Louvain approach with a refinement step that addresses poorly connected or internally disconnected communities \citep{Traag2019}. Its use of edge weights allows the strength of pairwise relationships to contribute to the partition. The resulting communities describe graph structure under the chosen objective function and resolution parameter; they do not establish that each community is a distinct outbreak.

When modularity is used as the objective, its resolution limit can favour merging small communities, even when they have distinct internal structure \citep{Fortunato2010}. Changing the resolution parameter can also change the size and membership of inferred communities. These properties make sensitivity to clustering choices part of the interpretation of genomic clusters. Chapters~\ref{ch:epilink} and~\ref{ch:genomic_networks} describe the use of Leiden in this thesis and the assessment of its resolution parameter.

\section{Cluster Transition Graph}\label{sec:ch2_transition_graph}

Recent-linkage clusters depend on how the graph and clusters are constructed. Each cluster groups genomes sampled within the same surveillance window. A cluster observed in one window cannot by itself show whether its members were represented in earlier clusters or will appear in later ones. Describing how cluster membership changes across windows requires linking clusters across adjacent surveillance windows. In this thesis, overlapping windows allow the same sampled genomes to appear in more than one window. Comparing membership yields a \emph{cluster-transition graph}. Its nodes represent clusters in particular windows, whereas compatibility-graph nodes represent sampled sequences. In descriptions of the transition graph, we use \emph{cluster} for this unit and \emph{link} for a connection based on shared membership between adjacent windows. Counts across windows therefore count separate cluster assignments, not distinct outbreaks. This graph describes how sampled membership continues or changes across windows.

Work on dynamic communities in network science provides terms for describing how groups emerge, persist, contract, expand, split, and merge \citep{Palla2007,Mucha2010,Saeed2025}. In temporal networks, combining links from different times into a single static graph can create paths that would be impossible if time order were retained, obscure causal ordering, and distort measures of spreading potential \citep{Holme2012,Kempe2002,Masuda2017}. The cluster-transition graph connects only clusters in adjacent windows, with edges directed from earlier to later windows. Every path therefore follows successively later surveillance windows. The shared samples defining those links do not establish the order of infection or transmission. Previous work tracked week-to-week changes in cluster size to identify possible superspreading \citep{dor2025}, using clusters defined by genetic similarity alone. These studies motivate examining the sudden appearance of large clusters and increases in sampled membership as surveillance signals \citep{GomezCarballa2021}.

Window overlap makes it possible to identify shared membership between clusters in neighbouring windows. It also introduces dependence: apparent persistence may arise partly because the same sampled genomes enter multiple windows. This differs from many social or contact networks, where the same individuals may be observed repeatedly at different times. In genomic surveillance, each sampled genome typically has one collection date. In this thesis, overlapping windows allow that genome to enter more than one clustering analysis. A link between an earlier and a later cluster therefore represents \emph{sampled continuity}, shaped by the window design, sampling process, and cluster definition. The transition graph summarises the surveillance record; it does not establish a sequence of transmission events. Chapter~\ref{ch:cluster_transition_screening} puts this interpretation into practice through a calibrated screening procedure. The following section explains how this screen identifies patterns compatible with superspreading from cluster membership and links across windows.

\section{Screening for Superspreading Signals}\label{sec:ch2_superspreading_signals}

The cluster-transition graph provides a way to identify patterns compatible with superspreading. As described in Section~\ref{sec:ch2_partial_observation}, genomic surveillance observes an incomplete and uneven sample of infections. Screening scores therefore need to be interpreted relative to comparison groups sampled under similar conditions. Comparing scores with a reference distribution is referred to here as \emph{calibration}. This section explains that comparison and why genomic surveillance alone cannot confirm superspreading events. The \emph{local burst} score combines cluster size relative to its sampled background with the proportion of members absent from directly linked earlier clusters, where those links are observed. The \emph{onward burden} score measures additional sampled infections in later linked clusters relative to the starting cluster's size.

\subsection{Statistical definition of a superspreading event}\label{sec:ch2_transmission_heterogeneity}

In infectious disease epidemics, some infected individuals cause many more secondary infections than others. This variation is described by the \emph{offspring distribution}: the probability distribution of the number of secondary infections, $Z$, caused by one infected person. In a fully susceptible population without control measures, its mean is the basic reproduction number, $\mathbb{E}[Z] = R_0$. A common model separates each case's expected number of secondary infections, $\nu$, from their actual count, $Z$. It assumes that $Z \mid \nu \sim \operatorname{Poisson}(\nu)$. If $\nu$ varies across cases according to a gamma distribution with mean $R_0$ and shape parameter $k$, then $Z$ follows a negative binomial distribution, with

\[
    \operatorname{Var}(Z) = R_0 + \frac{R_0^{2}}{k}.
\]

At a fixed $R_0$, smaller values of $k$ indicate greater variation between cases, with transmission concentrated among fewer infected individuals \citep{LloydSmith2005}. 

The parameter $k$ describes variation across cases in a particular population and context; it is not an individual characteristic or a label for a specific cluster. It reflects differences in infectiousness, contact patterns, network position, setting, and timing. In practice, interventions and behaviour can change the transmission pattern, while incomplete case detection can affect the estimated value of $k$. When describing transmission under interventions or existing immunity, a context-specific (or \emph{effective}) reproduction number $R$ should replace $R_0$.

The corresponding percentile definition operates at the individual level \citep{LloydSmith2005}. This approach estimates the effective reproduction number $R$ and constructs a $\operatorname{Poisson}(R)$ distribution. This represents the secondary-infection counts expected from random variation when all cases have the same expected number of secondary infections. An individual is associated with an \emph{$n$th-percentile superspreading event} if they cause more than $Z^{(n)}$ secondary infections, where $Z^{(n)}$ is the $n$th percentile of that distribution. Applying this definition requires three elements that are unavailable or difficult to infer from the Scottish surveillance record: a primary case to which the secondary infections can be attributed, an observed or reconstructed count of those infections, and a reference value of $R$ appropriate to the setting.

\subsection{Challenges in inferring superspreading events from genomic data}\label{sec:ch2_identifying_sses}

Superspreading has been documented for severe acute respiratory syndrome coronavirus 1 (SARS-CoV-1)\nomenclature{SARS-CoV-1}{severe acute respiratory syndrome coronavirus 1}, Middle East respiratory syndrome coronavirus (MERS-CoV)\nomenclature{MERS-CoV}{Middle East respiratory syndrome coronavirus}, Ebola virus, and SARS-CoV-2 \citep{Shen2004,Cho2016,Lau2017,Endo2020,McKee2024,Stein2011}. For SARS-CoV-2, epidemiological analyses found strong overdispersion in secondary-infection counts, with many infections causing no detected secondary cases and a minority accounting for a disproportionate share of onward transmission \citep{Endo2020,Adam2020,McKee2024}. Genomic analysis of the Boston epidemic demonstrated how transmission associated with particular events could contribute extensively to subsequent community spread \citep{Lemieux2021}. However, reconstructing the source and timing of such events remains difficult even with detailed epidemiological information. The Skagit County choir outbreak was initially interpreted as a large superspreading event from one rehearsal \citep{Hamner2020,Miller2020}, but subsequent analysis questioned whether symptom onset timing was consistent with all infections originating from that single event \citep{Axon2023}.

Concentrated transmission can produce groups of highly similar viral genomes sampled close together in time. Such patterns help identify transmission episodes but are not specific to superspreading. Similar patterns can arise through multiple transmission generations, multiple introductions of a common viral genotype, or changes in testing and sequencing intensity. Incomplete sampling can obscure the genomic signature of genuine superspreading events. Genomic and temporal proximity therefore indicate possible recent shared transmission history, not the number of secondary infections from a particular primary case. The contextual information needed to distinguish these explanations---infectiousness, contact opportunities, ventilation, timing of detection---is generally unavailable in routine national genomic surveillance \citep{Sun2021,Althouse2020,Stein2011}.

The clusters inferred in this thesis group sampled infections using compatibility scores based on genetic differences and collection dates within rolling windows. They are not groups defined solely by identical genomes or reconstructed transmission chains. A community may contain multiple transmission generations, non-identical but compatible genomes, or repeated introductions of similar viruses. Its membership depends on window boundaries, the recent-linkage scenarios specified in the compatibility model (see Section~\ref{sec:ch1_inferential_target}), sequencing coverage, genetic diversity between distinct outbreaks, and the community-detection procedure. Cluster size therefore denotes the number of sampled genomes assigned to a within-window community, not the number of secondary infections from one primary case. Without an identified primary case to which secondary infections can be attributed and a suitable reference reproduction number $R$, this count cannot be treated as an observation of $Z$ or assessed against an offspring-based superspreading threshold.

Cluster sizes may partly reflect variation in the number of onward infections, but they also depend on transmission over multiple generations, viral evolution, sampling, and cluster construction. Restricting samples to a time window can include only part of an ongoing transmission chain. The clustering procedure may split a chain into several groups or merge distinct chains. Sizes in different windows therefore need to be assessed against their sampling conditions, and changes in sampled cluster size do not by themselves establish growth in transmission. Large clusters are useful surveillance signals, but do not directly measure individual transmission or provide sufficient evidence of superspreading.

These limitations also prevent direct application of existing cluster-based methods.\ \citet{TranKiem2024} derived a likelihood for $R$ and $k$ from cluster size distributions, explicitly modelling the offspring distribution and sequencing fraction. Applying that approach here would require a model describing how transmission and sampling produce these particular compatibility communities.\ \citet{dor2025} tracked week-to-week changes in persistent genomic clusters defined by identical amino-acid substitution profiles. In this thesis's cluster-transition graph, each cluster is defined in one surveillance window. Differences between linked cluster sizes can therefore reflect changing window membership, community redefinition, splitting, merging, or sampling, rather than growth of a fixed transmission group.

\subsection{A calibrated screening approach}\label{sec:ch2_screening_approach}

The surveillance aim is to identify genomic patterns compatible with unusually concentrated transmission and prioritise them for further investigation. These patterns can be reviewed against the underlying samples and screening rules, without treating them as confirmed superspreading events or estimates of the offspring distribution. The procedure compares cluster size and the proportion of members absent from earlier linked clusters, where those links are observed, against relevant sampled backgrounds. It also assesses whether additional sampled infections in later linked clusters are unusually numerous relative to the starting cluster's size. These features describe the sampled genomic record; they do not establish new introductions, transmission from one cluster to another, or secondary-infection counts. Their value lies in identifying signals that merit closer examination, rather than explaining how transmission occurred from cluster size alone.

The screen assesses clusters against comparison groups defined using surveillance window, clade, and the availability of score components, with broader groups used when needed. Calibration asks whether a score is unusually high relative to these comparison groups, rather than simply large in absolute terms. A cluster in the transition graph is assigned candidate superspreading-compatible (CSC)\nomenclature{CSC}{candidate superspreading-compatible} status when its calibrated local burst score, onward burden score, or both meet the screening threshold. Chapter~\ref{ch:cluster_transition_screening} gives the full scoring and calibration procedure.

\section{Characterising Genomic Graphs and Clusters}\label{sec:ch2_characterisation}

The preceding sections describe compatibility graphs, recent-linkage clusters, local burst and onward burden scores, and CSC clusters. Characterisation examines the sampled populations, places, periods, lineages, and testing contexts represented in these results. This separates two questions: whether pairs of sampled infections are compatible with the selected recent-transmission scenarios, and which populations and surveillance conditions are represented in the resulting graphs, clusters, and screening outcomes. Table~\ref{tab:candidate_definition} defines these terms and distinguishes them from a superspreading event.

\begin{table}[htbp]
    \centering
    \caption[Definitions of graphs, clusters, screening outcomes, and superspreading]{Definitions of graphs, clusters, screening outcomes, and superspreading in this thesis.}\label{tab:candidate_definition}
    \begin{thesistablebody}{@{}P{0.28\textwidth}P{0.72\textwidth}@{}}
    \toprule
    \textbf{Term} & \textbf{Meaning in this thesis} \\
    \midrule
    Compatibility graph &
    A weighted graph for one surveillance window and lineage. Nodes represent sampled genomes, and edge weights describe compatibility with the selected recent-transmission scenarios. \\

    \addlinespace
    Recent-linkage cluster &
    A group of sampled genomes assigned to the same community within a compatibility graph, based on stronger compatibility within groups than between them. \\

    \addlinespace
    Cluster in the transition graph &
    A cluster from one surveillance window, connected to clusters in adjacent windows through shared membership. Incoming and outgoing links describe shared sampled membership between neighbouring windows. \\

    \addlinespace
    Screening scores &
    \emph{Local burst} combines cluster size relative to its sampled background with the proportion of members absent from directly linked earlier clusters, where those links are observed.\ \emph{Onward burden} measures additional sampled infections in later linked clusters relative to the starting cluster's size. \\

    \addlinespace
    Candidate superspreading-compatible (CSC) cluster &
    A cluster whose calibrated local burst score, onward burden score, or both are unusually high relative to comparable clusters. \\

    \addlinespace
    Superspreading event &
    At individual level, a case causes more secondary infections than the percentile threshold defined above. Broader event or setting usage describes an episode of unusually concentrated transmission supported by epidemiological evidence. \\
    \bottomrule
    \end{thesistablebody}
\end{table}

\subsection{Epidemiological context for characterisation}\label{sec:ch2_population_place_context}

The epidemiology of SARS-CoV-2 motivates examining the populations and places represented in genomic graphs and clusters. In the United Kingdom and elsewhere, risks of infection, severe disease, hospitalisation, and mortality varied significantly by socioeconomic deprivation, ethnicity, occupation, household composition, geography, and disability, reflecting and often amplifying pre-existing health inequalities \citep{Bambra2020,Niedzwiedz2020,Mathur2021,Whitehead2021,Marmot2020}. In Scotland, Public Health Scotland documented differences in COVID-19 mortality between places and deprivation groups, while linked-data research examined differences in hospital-admission risk and vaccine effectiveness between variants \citep{PublicHealthScotland2020,Sheikh2021}. These differences matter for genomic surveillance because the same social and geographic factors shape exposure opportunities, access to testing, sample collection, selection for sequencing, and linkage to other records.

The literature on superspreading also highlights settings and exposures that are only partly visible in routine surveillance data. Care homes, hospitals, shelters, prisons, workplaces, and large gatherings recur as settings where crowding, prolonged indoor exposure, patterns of close contact, and delayed detection can favour concentrated transmission. Where people have a higher risk of severe disease, an outbreak may also have more serious health consequences \citep{Althouse2020,Meredith2020,Hamilton2021,Aggarwal2022,Lemieux2021,McKee2024}. In Scotland, one study identified cases linked to the EURO 2020 European football championship through dedicated contact-tracing tags, illustrating how event context may require information not recoverable from residence, age, deprivation, or lineage alone \citep{Marsh2021}. Characterisation therefore asks whether the inferred clusters show patterns relevant to these settings and exposures. Residential and demographic variables provide indirect clues, but do not establish where transmission occurred.

\subsection{Questions supported by linked variables}\label{sec:ch2_contextual_variables}

The linked Scottish surveillance record supports descriptions of sampled infections, residential areas, and surveillance periods. At the \emph{sequence level}, records include demographic attributes, vaccination history, testing information, collection date, viral lineage and clade, and residential Data Zone. Data Zones are standard small-area geographical units used for official statistics in Scotland, each containing approximately 500 to 1,000 residents. At the \emph{area level}, they can be linked to local authorities, health boards, urban/rural classes, and the Scottish Index of Multiple Deprivation (SIMD)\nomenclature{SIMD}{Scottish Index of Multiple Deprivation} \citep{ScottishDataZones2022,ScottishSIMD2020,ScottishUrbanRural2020}. At the \emph{window level}, records can be interpreted against policy periods, recorded case counts, numbers of genomes sequenced, sequencing coverage, and lineage composition.

At the graph level, these variables support the composition and assortativity comparisons described above. Their interpretation depends on the sampled population and the meaning of each category; they do not establish contact between individuals.

For clusters in individual windows and their transition links, analysis can examine how size, residential spread, and shared membership across windows vary by lineage, period, sequencing context, or sampled population composition. It can also describe links from several earlier clusters into one later cluster, links from one cluster to several later clusters, or the absence of an observed link to a later cluster. The last pattern marks the end of an observed graph path, not necessarily the end of transmission.

For CSC clusters, the comparison is with other eligible sampled clusters. Analysis can examine which population groups, places, policy periods, clades, testing pathways, and sequencing contexts are over- or under-represented among selected clusters. Sequence-level composition compares sampled records within these cluster groups; a sequence can contribute in more than one window.

\section{Conclusion}\label{sec:ch2_conclusion}

The review motivates a pairwise measure that combines genetic differences with sampling dates, assesses explicitly defined recent-linkage scenarios, and can be used without fitting to labelled transmission pairs. Its interpretation must account for lineage background, surveillance processes, and model assumptions. Clusters and their changes over time require comparison with an appropriate sampled background, while characterisation establishes which populations those patterns represent. These requirements support screening for epidemiological investigation without treating genomic clusters as confirmed superspreading events.

Chapter~\ref{ch:epilink} addresses the pairwise requirement by developing and evaluating EpiLink.
